\documentclass{article}
\usepackage{graphicx} % Required for inserting images

\title{Ensayo sobre los beneficios de la programación en la astrofísica}
\author{María Salvador\\
       202506721}
\date{05 de febrero 2026}

\begin{document}

\maketitle

\section{Introducción}
 La astrofísica es un área de estudio de la física y la astronomía que aplica leyes físicas fundamentales, con el tiempo ha evolucionado, despúes de hacer observaciones por medio de telescopios en la actualidad la información generada por observatorios modernos es inmensa para el ojo humano.\\
 
 \section{Desarrollo}
 La programación permite el procesamiento automatizado y el análisis de grandes conjuntos de datos.
 Además, la simulación computacional es esencial para entender fenómenos inabarcables en un laboratorio, el software permite el análisis rápido de imágenes y datos de telescopios, lo que agiliza el descubrimiento de nuevos objetos celestes, permite recrear modelos físicos para comparar simulaciones teóricas con observaciones reales, lo que refuerza la comprensión física.\\

\section{Conclusión}
El futuro de la astrofísica está intrínsecamente ligado a la capacidad de innovar en el desarrollo de software y algoritmos. 


\end{document}
